\begin{table}[!htbp]
\centering
\footnotesize
\begin{threeparttable}
\caption{Placebo test of cross-country stress synchronization. Country-level stress indicators are independently reshuffled across calendar months, breaking common-shock structure while preserving country-level event frequencies.}
\label{tab:placebo}
\begin{tabular}{lrrrr}
\toprule
Statistic & Observed & Placebo mean & Placebo P95 & $p$-value \\
\midrule
Max avalanche size $S_{\max}$ & 11 & 3.47 & 4.00 & 0.000 \\
Mean avalanche size $\bar{S}$ & 3.50 & 1.32 & 1.40 & 0.000 \\
Share $P(S\geq 4)$ & 0.350 & 0.005 & 0.019 & 0.000 \\
Share $P(S\geq 6)$ & 0.250 & 0.000 & 0.000 & 0.000 \\
\bottomrule
\end{tabular}
\begin{tablenotes}\footnotesize
\item Bootstrap replications: $B=1000$, seed $=2026$.
\item Under the null of independent country-level stress, the maximum simultaneously-hit panel is at most $\bar{S}_{\max}^{\mathrm{pl}}\approx 3.5$. The observed value $S_{\max}=11$ is unattainable in the placebo distribution.
\item All four synchronization statistics reject the independence null at $p<0.001$, confirming that simultaneous LATAM stress is a non-trivial multi-country phenomenon.
\end{tablenotes}
\end{threeparttable}
\end{table}
